\section{Does the tradeoff hold with reported preferences?}\label{sec:field}
In the rental experiment, learning raised partner recovery and lowered placement relative to uniform exposure. Would the same two signs recur if the rankings came from people? We fixed a five-option comparison in each of four settings before computing its outcomes. We fitted each model to earlier cohorts. Features and targets differ by institution. In particular, school and university applicants already submitted their options. Our experiment curates those lists. It cannot recover preferences over programs they never listed. Table~\ref{tab:field-transport} gives every primary contrast against uniform.

\begin{table}[t]\centering\small
\caption{Learned five-option menus minus uniform, in percentage points. DA partner recovery is conditional on placement under full-list DA. The next column measures coverage of recorded offers or, in the last row, mutual yeses. Placement counts everyone on the proposing side. All placement contrasts are counterfactual. Each predictor uses the information available in its setting.}
\label{tab:field-transport}\begin{tabular}{@{}lrrr@{}}\toprule
Test cohort & DA partner shown & Recorded coverage & DA placement\\\midrule
NYC summer ordinary (simulated) & $+21.97$ & --- & $-3.08$\\
NYC winter ordinary (simulated) & $+21.90$ & --- & $-3.35$\\\midrule
Chile K--12, 2020 & $+1.81$ & $+1.80$ & $+0.28$\\
Chile K--12, 2021 & $+1.52$ & $+1.52$ & $+0.24$\\
Chile universities, 2026 & $+23.81$ & $+23.83$ & $+0.92$\\
University mentoring, third cohort & $+9.53$ & --- & $+0.67$\\
Columbia speed dating & $+1.37$ & $+6.71$ & $-5.64$\\\bottomrule
\end{tabular}\end{table}

\paragraph{National university admissions.}
Chile's regular admissions files let us distinguish an actual offer from a model's prediction target \citep{paes,paesrules}. The 2026 public cohort contains 189,399 applicants and 1,328,381 admissible regular applications across 2,150 published programs. Each application supplies a student rank and a program-specific weighted selection score. We take the advertised regular seats plus overbooking and prioritize applicants by score within each program. Full-list DA reproduces the exact program for 148,457 of 156,248 recorded regular recipients (95.01\%). Including the outside option, exact agreement is 93.59\% of all applicants. Reserved tracks, tie rules and administrative exclusions limit fidelity. We keep the 7,661 applicants with no eligible scored option in the placement denominator.

A regularized classifier fitted on 2025 full-list DA uses submitted rank, list length, and academic score relative to a program's previous cutoff. It retains the recorded offer for 98.05\% of historical recipients, against 74.22\% under uniform exposure. It also places about 1,741 more applicants than uniform. A simpler rule prioritizes programs where the applicant clears last year's minimum regular offer score, ordering each group by the applicant's reported rank. Against this rule, the trained model puts \emph{3,359 more recorded offers} on lists yet produces \emph{1,036 fewer placements} (157,562 versus 158,598). The coverage advantage also holds using full-list DA as the target (97.64\% versus 96.35\%).

A follow-up isolates access to submitted rank. We trained two copies of the same 4,514-parameter network on actual 2025 offers, then scored the same 2026 lists. Supplying the rank covers \emph{11,256 more actual offers} but places \emph{267 fewer applicants}. The rank-informed rule keeps 681 more full-list DA recipients placed while placing 948 fewer people who had been unmatched under full DA. This accounts for the entire net loss. Both models beat uniform on both outcomes. Appendix~\ref{app:paes} gives the protocol, including an eligibility check that uses no test-year outcome codes.

\paragraph{Two school admission cycles.}
The 2020 and 2021 K--12 tests cover 915,638 applicant records. A model trained on 2019 offers knows which courses a family listed and how many applicants and seats a course has. It does not see the family's ranking or the lottery draw. At five discretionary slots, the fitted rule improves both outcomes relative to uniform. Adding the reported preference positions to the model improves historical offer coverage by a further 0.32 and 0.33 points. It costs 525 and 519 placements in the respective years. Both models were specified before scoring. We examined their head-to-head contrast after the primary test. Most lists fit within five slots. Among people whose lists could not change, placement fell by 0.13 and 0.07 points under the rank-blind model relative to uniform. These are spillovers from changes to someone else's list. Full-list replay reproduces the exact recorded outcome, including nonassignment, for 95.24\% and 95.52\% (Appendix~\ref{app:sae}) \citep{saedata,correaschool}.

\paragraph{A directional test in two later cohorts.}
Before obtaining the next two application years, we fixed the 2019 model and the prediction that adding rank would raise actual-offer coverage \emph{and} reduce placement in each year. Both signs held in both cohorts. In 2022 the rank-informed menus covered 1,143 more recorded offers and placed 468 fewer of 570,891 applicants. In 2023 they covered 1,350 more and placed 517 fewer of 536,353. Of the 2,493 net additional historical offers covered, 2,244 were offers families explicitly accepted. The placement mechanism differs from the university example. Across the two school cohorts, the informed rule retained 227 fewer full-list recipients and admitted 758 fewer previously unmatched applicants. When only applicants outside the ministry's priority-designated group received the new rule, 252 designated pupils lost a placement and 127 gained one despite keeping their own lists. A reserve-first allocation sensitivity preserved both overall signs. Appendix~\ref{app:sae} gives the denominators, gross flows, and the distinction between initial acceptance and attendance.

\paragraph{Does an offer stand in for a choice the applicant actually made?}
We linked 136,751 university applicants to a 2026 registration record. Of their enrolled programs, 5,271 never appeared on an eligible regular list and could not be recovered by pruning one. Counting these as misses, the original rank-informed offer predictor covered 129,668 enrolled destinations, compared with 118,898 for its rank-blind counterpart. The difference of 10,770 includes 1,305 among students who registered somewhere other than their initial regular offer and 954 among those without such an offer. We observe the registrations under the historical system. Whether anyone would register following a different placement is a separate question.

To check whether this was peculiar to training on offers, we refit both networks using 2025 enrollment destinations. Against the otherwise identical enrollment-trained model without rank, adding rank covered 10,891 more actual 2026 registrations and placed 141 fewer applicants under DA. This was a follow-up on the same test cohort. Appendix~\ref{app:enroll} gives the target construction, counts and full-list recipient changes.

\paragraph{When both sides report a preference.}
The held-out mentoring cohort of \citet{haashall2019} contains 29 mentees. A fitted rule using both sides' preference tiers improves both metrics against uniform. Compared with a model that sees only the graph and its degrees, it recovers more partners (97.95\% versus 94.00\%) but places fewer mentees (91.89\% versus 92.47\%). The speed-dating comparison comes closest to learning before preferences are revealed. In eight later sessions with 252 participants, a pre-event model retains a larger fraction of actual mutual yeses (41.43\% versus 34.72\%). Under one-partner DA, its two point estimates have the rental signs, but both date-clustered intervals span zero. A simple age-and-interest score improves both DA point estimates. The event allowed multiple mutually accepted contacts. DA models an alternative one-contact institution \citep{fisman2006dating,bhargava2014data}.

\paragraph{What limits the number of placements?}
Let $\nu(D)$ be the greatest number a display can place given its capacities, ignoring stability. Write $g(D)=\nu(D)-|\da(D)|$. The placement difference is exactly $\Delta\nu-\Delta g$. In ordinary NYC summer markets, contextual ranking reduces visible capacity by 5.19 percentage points and reduces the DA shortfall by 2.11, giving the observed 3.08-point loss. Winter gives $-5.60-(-2.25)=-3.35$. In the 2026 university comparison, the trained rule's display has room for 906 fewer placements than the prior-cutoff display (177,406 versus 178,312), and its DA shortfall grows by 130. These account for all 1,036 fewer placements. In mentoring, both fitted displays can place all 29 mentees. The difference comes entirely from the assignments DA selects. Maximum-cardinality assignments provide an upper bound and can violate priorities.
